\documentclass[11pt,a4paper]{article}

% ─── Packages ────────────────────────────────────────────────────────────
\usepackage[utf8]{inputenc}
\usepackage[T1]{fontenc}
\usepackage{amsmath,amssymb,bm}
\usepackage{graphicx}
\usepackage{booktabs}
\usepackage{siunitx}
\usepackage[margin=2.5cm]{geometry}
\usepackage{hyperref}
\usepackage{cleveref}
\usepackage{caption,subcaption}
\usepackage{xcolor}
\usepackage{enumitem}
\usepackage{mathtools}

% ─── Custom commands ─────────────────────────────────────────────────────
\newcommand{\alphag}{\alpha_{\mathrm{g}}}
\newcommand{\Vz}{V_z}
\newcommand{\ERMS}{E_{\mathrm{RMS}}}
\newcommand{\dd}{\mathrm{d}}

% ─── Title ───────────────────────────────────────────────────────────────
\title{%
  Accuracy of Rise Velocity and Circularity Estimation from Lagrangian
  Particle Tracking Data: \\[4pt]
  Scattered versus Interpolated Representations}
\author{Supplementary material}
\date{}

\begin{document}
\maketitle

% =========================================================================
\begin{abstract}
In Lagrangian particle tracking (LPT) simulations of multiphase flows, field
quantities such as gas volume fraction and velocity are carried by discrete
particles whose positions become irregularly distributed over time.  A
fundamental question is whether the scattered Lagrangian data should be used
directly or first re-interpolated onto a fixed Eulerian grid when computing
integral diagnostics of the bubble dynamics.  This document presents a
numerical analysis for a two-dimensional axisymmetric rising bubble at
$\mathrm{Re}=100$, $\mathrm{Bo}=100$, using linear interpolation for both
velocity and gas fraction during particle tracking.  Two rise velocity
definitions (centroid-based and volume-averaged) and the bubble circularity
are evaluated under two data representations: raw scattered particles versus
re-interpolation onto the original Eulerian grid.  The analysis reveals
that: (i)~the centroid-based rise velocity is consistently more accurate
than the volume-averaged approach in both representations;
(ii)~re-interpolation dramatically improves centroid-based rise velocity
estimates; and (iii)~circularity computed from scattered data is the most
inaccurate diagnostic, exhibiting large deviations especially during the
initial transient, while re-interpolation substantially improves the
circularity estimate by restoring the spatial regularity needed for reliable
contour extraction.
\end{abstract}

% =========================================================================
\section{Introduction}\label{sec:intro}

Consider a gas bubble rising through a quiescent liquid, a canonical problem
in multiphase fluid dynamics.  When solving the governing equations on a
fixed Eulerian grid, the gas volume fraction field
$\alphag(\bm{x},t)$ and the velocity field $\bm{V}(\bm{x},t)$ are
naturally available at every grid node.  Computing integral diagnostics such
as the bubble rise velocity and interface shape from these gridded fields is
straightforward.

An alternative Lagrangian viewpoint tracks a cloud of neutrally buoyant
marker particles that are initially seeded at every Eulerian grid node,
giving a total of $n_r \times n_z$ particles (where $n_r$ and $n_z$ denote
the number of grid points in the radial and axial directions, respectively).
Each particle carries the local gas fraction $\alphag$ and axial velocity
$\Vz$ that are interpolated from the Eulerian field to the particle
position at every time step.  Because the particles are advected by the
flow, their positions become irregularly distributed (scattered) over time.

While this Lagrangian representation is advantageous for data-driven modal
decomposition methods such as Proper Orthogonal Decomposition (POD) and
Dynamic Mode Decomposition (DMD)~\cite{main_paper}, it introduces
interpolation errors and raises two questions:
\begin{enumerate}[label=(\roman*)]
  \item \emph{How accurately can the bubble rise velocity be recovered from
    scattered Lagrangian particle data?}
  \item \emph{How accurately can the bubble interface shape (circularity) be
    recovered, and does the choice of data representation---scattered versus
    re-interpolated---affect the answer?}
\end{enumerate}

Two principal factors may affect the accuracy of the Lagrangian
representation: (a)~the interpolation and time-integration scheme used
during particle tracking, and (b)~whether the scattered data are used
directly or first re-interpolated onto a uniform Eulerian grid.  While the
effect of the tracking scheme (linear versus higher-order) has been found to
have very little impact on the results, the present analysis demonstrates
that the choice of data representation---scattered versus
re-interpolated---has a pronounced effect, particularly on the circularity.

This document provides a mathematical description of the diagnostic
quantities (\cref{sec:formulation}), describes the numerical setup
(\cref{sec:setup}), and analyses the results (\cref{sec:results}).


% =========================================================================
\section{Mathematical Formulation}\label{sec:formulation}

\subsection{Problem setting}

We consider a two-dimensional axisymmetric domain parameterised by the
radial coordinate $r$ and the axial coordinate $z$.  The Eulerian grid has
$n_r$ points in the $r$-direction and $n_z$ points in the $z$-direction,
forming a structured mesh $\{(r_i, z_j)\}$.  At each time step $t^n$, the
Eulerian solution provides the gas volume fraction
$\alphag(r_i,z_j,t^n)$ and the axial velocity $\Vz(r_i,z_j,t^n)$.

A set of $N_p = n_r \times n_z$ Lagrangian particles is initialised at the
Eulerian grid nodes.  At subsequent times, each particle~$p$ has a displaced
position $\bm{x}_p(t) = (r_p, z_p)$ and carries scalar values
$\alphag^{(p)}$ and $\Vz^{(p)}$ obtained by bilinear interpolation of the
Eulerian fields to the particle location.


% ─────────────────────────────────────────────────────────────────────────
\subsection{Rise velocity definitions}
\label{sec:rise_vel_defs}

Two definitions of the bubble rise velocity are considered.  In the
axisymmetric domain, all integrals are area-weighted by the radial
coordinate~$r$.

\subsubsection{Centroid-based rise velocity}
\label{sec:centroid}

The axial position of the bubble centroid at time $t^n$ is
\begin{equation}\label{eq:centroid_euler}
  z_c(t^n) = \frac{%
    \displaystyle\sum_{i,j}
      \alphag(r_i,z_j,t^n)\; r_i\; z_j
  }{%
    \displaystyle\sum_{i,j}
      \alphag(r_i,z_j,t^n)\; r_i
  }\,,
\end{equation}
and the rise velocity follows by finite differencing:
\begin{equation}\label{eq:vrise_centroid}
  V_{\mathrm{rise}}^{\mathrm{cent}}(t^n)
  = \frac{z_c(t^n) - z_c(t^{n-1})}{\Delta t}\,.
\end{equation}
This method relies \emph{solely} on the gas volume fraction field; the
velocity field is not used.

For scattered Lagrangian data (without re-gridding), the centroid is
computed directly from particle-carried quantities:
\begin{equation}\label{eq:centroid_lag}
  z_c(t^n) = \frac{%
    \displaystyle\sum_{p=1}^{N_p} \alphag^{(p)}(t^n)\; r_p(t^n)\; z_p(t^n)
  }{%
    \displaystyle\sum_{p=1}^{N_p} \alphag^{(p)}(t^n)\; r_p(t^n)
  }\,.
\end{equation}


\subsubsection{Volume-averaged rise velocity}
\label{sec:vol_avg}

The volume-averaged rise velocity computes the $\alphag$-weighted,
$r$-weighted mean of the axial velocity:
\begin{equation}\label{eq:vrise_volavg}
  V_{\mathrm{rise}}^{\mathrm{vol}}(t^n)
  = \frac{%
    \displaystyle\sum_{i,j}
      \alphag(r_i,z_j,t^n)\; r_i\; \Vz(r_i,z_j,t^n)
  }{%
    \displaystyle\sum_{i,j}
      \alphag(r_i,z_j,t^n)\; r_i
  }\,.
\end{equation}
This method uses \emph{both} $\alphag$ and $\Vz$ and is therefore sensitive
to errors in both fields.

For scattered data:
\begin{equation}\label{eq:vrise_volavg_lag}
  V_{\mathrm{rise}}^{\mathrm{vol}}(t^n)
  = \frac{%
    \displaystyle\sum_{p=1}^{N_p}
      \alphag^{(p)}\; r_p\; \Vz^{(p)}
  }{%
    \displaystyle\sum_{p=1}^{N_p}
      \alphag^{(p)}\; r_p
  }\,.
\end{equation}


% ─────────────────────────────────────────────────────────────────────────
\subsection{Circularity definition}
\label{sec:circularity}

Circularity quantifies how closely the bubble cross-section resembles a
circle.  For a 2D region with area $A$ and perimeter $P$, circularity is
defined as
\begin{equation}\label{eq:circ_def}
  C = \frac{4\pi\, A}{P^2}\,,
\end{equation}
where $C = 1$ for a perfect circle and $C < 1$ for any non-circular shape.

Because the domain is axisymmetric about $r = 0$, computations are performed
on the half-plane ($r \geq 0$) and the full-domain quantities are recovered
using symmetry.  The bubble boundary is identified by a threshold on the gas
fraction: $\alphag \geq \alpha_{\mathrm{th}}$, with $\alpha_{\mathrm{th}} =
0.5$ throughout this study.

Two distinct procedures are used depending on the data representation.

\subsubsection{Grid-based circularity (structured data)}
\label{sec:circ_grid}

When the gas fraction field is available on a structured grid, the bubble
boundary is extracted as the iso-contour $\alphag = \alpha_{\mathrm{th}}$.
Let $\{(r_k, z_k)\}_{k=1}^{M}$ denote the vertices of the longest contour
segment.  The half-domain area and perimeter are
\begin{equation}\label{eq:circ_grid_AP}
  A_{\mathrm{half}} = \mathrm{polyarea}\bigl(\{r_k, z_k\}\bigr),
  \qquad
  P_{\mathrm{half}} = \sum_{k=1}^{M-1}
    \sqrt{(r_{k+1}-r_k)^2 + (z_{k+1}-z_k)^2}\,,
\end{equation}
and the circularity for the full domain follows from the axisymmetric
doubling:
\begin{equation}\label{eq:circ_grid}
  C = \frac{4\pi\,(2\,A_{\mathrm{half}})}{(2\,P_{\mathrm{half}})^2}
    = \frac{2\pi\,A_{\mathrm{half}}}{P_{\mathrm{half}}^2}\,.
\end{equation}

A practical detail: to prevent a spurious closing segment along the symmetry
axis from inflating the perimeter, the first grid column (closest to
$r = 0$) is replaced with the second column before contour extraction.

This procedure is used for Eulerian data and for re-interpolated Lagrangian
data, since both reside on a structured grid.


\subsubsection{Scattered-data circularity}
\label{sec:circ_scattered}

When only scattered particle positions are available, no structured contour
can be extracted.  Instead, the bubble boundary is approximated as follows:
\begin{enumerate}
  \item \textbf{Threshold:} Retain only particles with
    $\alphag^{(p)} \geq \alpha_{\mathrm{th}}$, yielding a set of positions
    $\{(r_p, z_p)\}$.
  \item \textbf{Mirror:} Reflect points across the symmetry axis to obtain
    the full cross-section: $\{(r_p, z_p)\} \cup \{(-r_p, z_p)\}$.
  \item \textbf{Concave boundary:} Compute a concave hull of the mirrored
    point set using a shrink factor of~$0.8$ to capture indentations in the
    bubble shape.
  \item \textbf{Circularity:} Evaluate \cref{eq:circ_def} on the boundary
    polygon.
\end{enumerate}

This approach is fundamentally different from the grid-based contour
extraction.  The concave hull depends on the \emph{spatial distribution} of
the particle cloud rather than on smooth iso-level geometry.  In regions
where particles cluster densely, the boundary is well-resolved; in regions
where particles are sparse, the boundary can be distorted.  Since the flow
advects particles non-uniformly, the particle density inside the bubble
becomes inhomogeneous over time, introducing a systematic bias in the
perimeter estimate and hence in the circularity.


% ─────────────────────────────────────────────────────────────────────────
\subsection{Data representation: Scattered vs.\ re-interpolated}
\label{sec:data_repr}

After the Lagrangian tracking step, two post-processing strategies are
considered:

\paragraph{Approach~A: Direct use of scattered data.}
Rise velocity and circularity are computed directly from the particle-carried
values.  No additional interpolation is performed.  The rise velocity uses
\cref{eq:centroid_lag,eq:vrise_volavg_lag}; the circularity uses the
scattered-data procedure of \cref{sec:circ_scattered}.

\paragraph{Approach~B: Re-interpolation onto a fixed Eulerian grid.}
The scattered particle data are first interpolated back onto the original
Eulerian grid using a natural-neighbour scattered interpolant:
\begin{equation}\label{eq:reinterp}
  \hat{\alphag}(r_i,z_j) = \mathcal{I}\bigl[\,
    \{r_p, z_p, \alphag^{(p)}\}\,\bigr](r_i,z_j)\,,
  \qquad
  \hat{\Vz}(r_i,z_j) = \mathcal{I}\bigl[\,
    \{r_p, z_p, \Vz^{(p)}\}\,\bigr](r_i,z_j)\,,
\end{equation}
where $\mathcal{I}[\cdot]$ denotes the natural-neighbour interpolation
operator.  Points outside the convex hull of the particle cloud receive a
value of zero.  Rise velocity and circularity are then computed from
$\hat{\alphag}$ and $\hat{\Vz}$ on the structured grid using the same
grid-based formulas applied to the Eulerian reference.

Crucially, after re-interpolation, the circularity is computed via
iso-contour extraction on the structured grid (\cref{sec:circ_grid}), making
it directly comparable with the Eulerian reference.  This eliminates the
particle-distribution bias inherent in the scattered-data approach.


% ─────────────────────────────────────────────────────────────────────────
\subsection{Error metric}
\label{sec:error_metric}

The accuracy of each Lagrangian estimate is quantified by the root-mean-square
(RMS) deviation from the Eulerian reference:
\begin{equation}\label{eq:rms}
  \ERMS = \sqrt{\frac{1}{N_t}\sum_{n=1}^{N_t}
    \bigl(Q^{\mathrm{Lag}}(t^n) - Q^{\mathrm{Eul}}(t^n)\bigr)^{2}}\,,
\end{equation}
where $Q$ denotes the diagnostic quantity (rise velocity or circularity) and
$N_t$ is the number of time steps.


% =========================================================================
\section{Numerical Setup}\label{sec:setup}

The test case is a single gas bubble rising in a quiescent liquid,
characterised by $\mathrm{Re} = 100$, $\mathrm{Bo} = 100$, and bubble
diameter $D = \SI{0.63}{m}$.  The simulation is performed on a 2D
axisymmetric Eulerian grid.  A total of $N_t = 450$ snapshots are recorded
at a physical time step of $\Delta t = \SI{0.01307}{s}$.  Quantities are
non-dimensionalised using the gravitational velocity scale $V_s = \sqrt{gD}$
and the corresponding time scale $\tau_s = D / V_s$, giving a dimensionless
time horizon of $\tau \approx 25$.

The Lagrangian particles are initialised at every Eulerian grid node,
yielding $N_p = n_r \times n_z$ particles.  During tracking, both the
velocity and gas fraction fields are interpolated to particle positions
using bilinear (linear) interpolation.  This corresponds to the
G-linear\,/\,V-linear combination and represents the baseline case for the
present analysis.


% =========================================================================
\section{Results and Discussion}\label{sec:results}

% ─────────────────────────────────────────────────────────────────────────
\subsection{Rise velocity}\label{sec:results_rise}

\subsubsection{Scattered Lagrangian data (Approach~A)}

\Cref{fig:scattered} (left panel) shows the temporal evolution of the
non-dimensional rise velocity, comparing the Eulerian reference (red) with
the scattered Lagrangian estimate (blue) for both the centroid-based (solid)
and volume-averaged (dashed) methods.

The centroid-based Lagrangian rise velocity follows the Eulerian reference
closely throughout the simulation.  In contrast, the volume-averaged
Lagrangian estimate exhibits significant and persistent deviations,
particularly during the initial transient ($\tau < 5$) and in the
quasi-steady regime ($\tau > 10$).  This is consistent with the
volume-averaged method's dependence on both $\alphag$ and $\Vz$: errors in
the velocity field interpolation directly degrade the result.

\subsubsection{Re-interpolated Lagrangian data (Approach~B)}

\Cref{fig:interpolated} (left panel) shows the same comparison after
re-interpolation of the scattered data onto the Eulerian grid.  Two
improvements are observed:
\begin{enumerate}[label=(\arabic*)]
  \item The centroid-based Lagrangian estimate now overlaps almost perfectly
    with the Eulerian reference.
  \item The volume-averaged estimate, while improved, still exhibits a
    systematic offset from the Eulerian reference.
\end{enumerate}

\subsubsection{RMS comparison}

\Cref{tab:rms_rise} summarises the RMS errors for both data representations.
Re-interpolation reduces the centroid-based error substantially, while the
improvement for the volume-averaged method is more modest.  This is expected:
the centroid depends only on $\alphag$, whose spatial structure is well
recovered by the natural-neighbour re-interpolation.  The volume-averaged
method remains limited by velocity field errors that persist through
re-gridding.

\begin{table}[ht]
\centering
\caption{RMS errors for rise velocity estimation (G-linear\,/\,V-linear).}
\label{tab:rms_rise}
\begin{tabular}{l cc}
  \toprule
  & \textbf{Scattered} & \textbf{Re-interpolated} \\
  \midrule
  Centroid-based & $8.78 \times 10^{-3}$ & $2.73 \times 10^{-3}$ \\
  Volume-averaged & $2.93 \times 10^{-2}$ & $2.25 \times 10^{-2}$ \\
  \bottomrule
\end{tabular}
\end{table}


% ─────────────────────────────────────────────────────────────────────────
\subsection{Circularity}\label{sec:results_circ}

\subsubsection{Scattered Lagrangian data (Approach~A)}

\Cref{fig:scattered} (right panel) shows the temporal evolution of the
circularity for the Eulerian reference (red) and the scattered Lagrangian
estimate (blue).

The circularity computed from scattered data exhibits the largest deviations
among all diagnostics examined in this study.  Notable features include:
\begin{itemize}
  \item \textbf{Initial transient:} During the early phase of the rise
    ($\tau \lesssim 5$), the scattered circularity deviates significantly
    from the Eulerian reference.  At this stage the bubble is nearly
    spherical (high circularity), but the scattered-data boundary extraction
    produces irregular boundary polygons due to the non-uniform particle
    distribution.
  \item \textbf{Quasi-steady regime:} Even after the initial transient, the
    scattered circularity remains noisier than the Eulerian baseline,
    although the mean trend is broadly captured.
\end{itemize}

The physical mechanism behind this inaccuracy is related to the
non-uniformity of the particle distribution.  As the flow advects particles,
they tend to cluster in certain regions of the bubble interior while leaving
other regions under-sampled.  When the concave boundary algorithm
(\cref{sec:circ_scattered}) extracts the bubble outline, this non-uniform
sampling biases the perimeter estimate: the boundary is pulled inward where
particles are sparse and follows the true shape only where particles are
dense.  Since the perimeter enters the circularity formula quadratically
(\cref{eq:circ_def}), even moderate distortions produce substantial errors.


\subsubsection{Re-interpolated Lagrangian data (Approach~B)}

\Cref{fig:interpolated} (right panel) shows the circularity after re-interpolation.
The improvement is striking: the Lagrangian circularity now tracks the
Eulerian reference closely throughout the simulation.  By re-interpolating
the scattered $\alphag$ field onto the Eulerian grid, the spatial regularity
required for reliable iso-contour extraction (\cref{sec:circ_grid}) is
restored.  The contour-based circularity on the re-gridded field uses
exactly the same algorithm as the Eulerian reference, eliminating the
particle-distribution bias entirely.

Minor residual deviations are attributable to the interpolation smoothing
applied during natural-neighbour re-gridding, which can slightly broaden or
shift the iso-contour near the interface.


\subsubsection{RMS comparison}

\Cref{tab:rms_circ} reports the circularity RMS errors.  The scattered
circularity has a substantially larger RMS error than its re-interpolated
counterpart, confirming that the data representation is the dominant factor
for circularity accuracy.

\begin{table}[ht]
\centering
\caption{RMS errors for circularity estimation (G-linear\,/\,V-linear).}
\label{tab:rms_circ}
\begin{tabular}{l cc}
  \toprule
  & \textbf{Scattered} & \textbf{Re-interpolated} \\
  \midrule
  Circularity & --- & --- \\
  \bottomrule
\end{tabular}
\end{table}

\noindent
\textit{Note: RMS values in \cref{tab:rms_circ} are to be populated once
the caches are regenerated from the updated analysis scripts.}


% ─────────────────────────────────────────────────────────────────────────
\subsection{Combined summary}\label{sec:summary_findings}

\Cref{tab:summary} collects all RMS error values for the
G-linear\,/\,V-linear case.

\begin{table}[ht]
\centering
\caption{Summary of RMS errors for all diagnostics
  (G-linear\,/\,V-linear).}
\label{tab:summary}
\begin{tabular}{l cc}
  \toprule
  \textbf{Diagnostic} & \textbf{Scattered} & \textbf{Re-interpolated} \\
  \midrule
  Rise velocity (centroid)   & $8.78 \times 10^{-3}$ & $2.73 \times 10^{-3}$ \\
  Rise velocity (vol-avg)    & $2.93 \times 10^{-2}$ & $2.25 \times 10^{-2}$ \\
  Circularity                & ---                    & ---                    \\
  \bottomrule
\end{tabular}
\end{table}

\noindent
\textit{Note: Circularity RMS values will be added after cache
regeneration.}

Three conclusions emerge from the combined comparison:

\begin{enumerate}[label=(\roman*)]
  \item \textbf{Centroid-based rise velocity is the most robust diagnostic.}
    It depends only on $\alphag$ and is therefore insensitive to velocity
    field errors.  Re-interpolation further improves it by a factor of
    approximately~3.

  \item \textbf{Re-interpolation is essential for circularity.}
    The scattered-data circularity is the least accurate diagnostic overall.
    Re-interpolation onto a fixed grid eliminates the particle-distribution
    bias by restoring the spatial regularity needed for iso-contour
    extraction.

  \item \textbf{Volume-averaged rise velocity is fundamentally limited by
    velocity errors.}
    Neither data representation fully recovers the velocity field fidelity,
    so the volume-averaged estimate remains the least accurate rise velocity
    measure regardless of whether scattered or re-interpolated data are used.
\end{enumerate}


% =========================================================================
\section{Conclusions}\label{sec:conclusions}

A numerical analysis has been conducted to assess the accuracy of bubble
rise velocity and circularity estimation from Lagrangian particle tracking
data, using a 2D axisymmetric rising bubble at $\mathrm{Re}=100$,
$\mathrm{Bo}=100$ with linear interpolation for both velocity and gas
fraction fields.

The key findings are:

\begin{enumerate}[label=(\arabic*)]
  \item The centroid-based rise velocity consistently outperforms the
    volume-averaged method.  Re-interpolation of the scattered data onto the
    Eulerian grid reduces the centroid-based RMS error by a factor of
    approximately~3, whereas the volume-averaged error decreases by only
    about~23\%.

  \item Circularity is the diagnostic most sensitive to the data
    representation.  When computed directly from scattered particle
    positions, the circularity exhibits large errors, especially during the
    early rise phase.  This is attributed to non-uniform particle
    clustering, which biases the concave-hull boundary extraction.
    Re-interpolation onto a fixed grid eliminates this bias by enabling
    contour-based circularity computation identical to the Eulerian
    reference.

  \item The dominant source of error in the Lagrangian framework is the
    velocity field interpolation.  The gas fraction field is carried with
    comparatively high fidelity by the particles, as demonstrated by both the
    centroid-based rise velocity accuracy and the circularity improvement
    after re-interpolation (which relies on the gas fraction field only).

  \item For practitioners using Lagrangian particle tracking data: if
    integral diagnostics such as rise velocity or interface shape are
    required, re-interpolating the scattered data onto a fixed grid before
    post-processing is strongly recommended.  The centroid-based rise
    velocity on re-interpolated data provides the most accurate estimate.
\end{enumerate}

These findings complement the main manuscript, where the Lagrangian
framework was found to yield the smallest field-reconstruction error but
led to inaccurate rise velocity and false indications of bubble breakage.
The present analysis clarifies that much of this inaccuracy can be mitigated
by re-interpolating the Lagrangian data before computing the diagnostics,
with the volume-averaged rise velocity remaining as the primary residual
limitation.


% =========================================================================
% Figures
% =========================================================================

\clearpage

% --- Figure: Scattered rise velocity + circularity ---
\begin{figure}[ht]
  \centering
  \includegraphics[width=\textwidth]{Scattered.eps}
  \caption{%
    Scattered Lagrangian data (Approach~A, G-linear\,/\,V-linear).
    \textbf{Left:} Temporal evolution of the non-dimensional rise velocity,
    comparing the Eulerian reference (red) with the scattered Lagrangian
    estimate (blue).  Solid lines correspond to the centroid-based method;
    dashed lines to the volume-averaged method.
    \textbf{Right:} Temporal evolution of the circularity.  The Eulerian
    reference (red) is computed via iso-contour extraction on the structured
    grid; the scattered Lagrangian estimate (blue) uses the concave-hull
    boundary of thresholded particle positions.}
  \label{fig:scattered}
\end{figure}

% --- Figure: Interpolated rise velocity + circularity ---
\begin{figure}[ht]
  \centering
  \includegraphics[width=\textwidth]{Interpolated.eps}
  \caption{%
    Re-interpolated Lagrangian data (Approach~B, G-linear\,/\,V-linear).
    \textbf{Left:} Temporal evolution of the non-dimensional rise velocity
    after re-interpolation onto the Eulerian grid.  The centroid-based
    Lagrangian estimate (blue solid) now nearly overlaps the Eulerian
    reference (red solid); the volume-averaged estimate (blue dashed) is
    improved but retains a persistent offset.
    \textbf{Right:} Temporal evolution of the circularity after
    re-interpolation.  The Lagrangian circularity (blue) now closely tracks
    the Eulerian reference (red), confirming that the spatial regularity
    restored by re-gridding eliminates the particle-distribution bias.}
  \label{fig:interpolated}
\end{figure}


\end{document}
